\documentclass{beamer}

\usepackage{beamerthemeSingapore}
\usepackage[brazil]{babel}
\usepackage[T1]{fontenc}
\usepackage[utf8]{inputenc}
\usepackage{graphicx}
\usepackage{url}
\usepackage{xcolor}
\usepackage{booktabs}
\usepackage{ragged2e}
\usepackage{etoolbox}
\usepackage{minted}
\usepackage{tikz}
\usetikzlibrary{positioning, arrows.meta, shapes.geometric}

\newcommand{\tcb}[1]{\textcolor{blue}{#1}}
\newcommand{\tcr}[1]{\textcolor{red}{#1}}
\newcommand{\tcg}[1]{\textcolor{green!50!black}{#1}}

\setminted{
  fontsize=\footnotesize,
  breaklines=true,
  breakanywhere=true,
  frame=lines,
  framesep=2mm,
  baselinestretch=1.05,
  tabsize=2
}

\apptocmd{\frame}{}{\justifying}{}

\setbeamertemplate{navigation symbols}{}

\title{03 - Protocolo HTTP}
\author{Prof. Alex Kutzke}
\date{Agosto de 2026}

\begin{document}

% ==========================================
\begin{frame}
  \begin{center}
    \large{Setor de Educação Profissional e Tecnológica}\\
    Tecnologia em Análise e Desenvolvimento de Sistemas\\
    \vspace{1cm}
    \small{DS122 - Desenvolvimento de Aplicações Web 1}
  \end{center}
  \titlepage
\end{frame}

% ==========================================
\begin{frame}{Objetivos da aula}
  Ao final desta aula você deve ser capaz de:
  \begin{enumerate}
    \item Descrever a estrutura de uma mensagem HTTP byte a byte;
    \item Explicar o que significa o protocolo ser \textit{sem estado} e as
          consequências disso;
    \item Identificar os três lugares por onde o cliente envia parâmetros:
          caminho, consulta e corpo;
    \item Escolher o método HTTP correto, justificando com segurança e
          idempotência;
    \item Interpretar códigos de status e cabeçalhos;
    \item Construir requisições manualmente com \texttt{curl};
    \item Atuar como servidor, respondendo à mão a uma requisição do navegador.
  \end{enumerate}
\end{frame}

% ==========================================
\begin{frame}{Roteiro}
  \tableofcontents
\end{frame}

% ==========================================
\section{Fundamentos}
% ==========================================

\begin{frame}{Modelo cliente-servidor}
  \begin{center}
    \begin{tikzpicture}[
        box/.style={draw, rounded corners, minimum width=3cm,
                    minimum height=1.6cm, align=center, thick},
        every node/.style={font=\small}
      ]
      \node[box, fill=blue!8]  (c) at (0,0)  {\textbf{Cliente}\\navegador, \texttt{curl}};
      \node[box, fill=green!8] (s) at (7,0)  {\textbf{Servidor}\\nginx, Apache};

      \draw[-{Stealth[length=3mm]}, thick, blue]
        ([yshift=4mm]c.east) -- node[above, font=\footnotesize]{requisição}
        ([yshift=4mm]s.west);
      \draw[-{Stealth[length=3mm]}, thick, green!50!black]
        ([yshift=-4mm]s.west) -- node[below, font=\footnotesize]{resposta}
        ([yshift=-4mm]c.east);
    \end{tikzpicture}
  \end{center}

  \begin{itemize}
    \item \textbf{Cliente}: inicia a comunicação. Envia uma requisição e aguarda.
    \item \textbf{Servidor}: nunca inicia. Fica escutando e responde.
  \end{itemize}

  \vspace{0.2cm}
  \footnotesize No modelo \textbf{desktop} existe um ator só: o programa roda na
  máquina do usuário, lê o teclado, escreve na tela, acessa arquivos locais.
\end{frame}

\begin{frame}{O servidor nunca inicia a comunicação}
  \begin{block}{}
    O servidor não tem como avisar espontaneamente o navegador de que algo
    mudou. Toda informação que chega ao cliente chega \textbf{porque o cliente
    pediu}.
  \end{block}

  \vspace{0.5cm}
  Uma página que precisa exibir dados novos sem interação do usuário só tem uma
  saída dentro do HTTP: perguntar de novo, de tempos em tempos. Contornar essa
  limitação exige tecnologias construídas por cima do protocolo, como WebSocket
  e \textit{Server-Sent Events}, fora do escopo desta disciplina.
\end{frame}

\begin{frame}{O que é um protocolo}
  Duas máquinas conectadas conseguem trocar bytes. Porém, isso geralmente não é
  suficiente: elas precisam concordar sobre o \textbf{significado} desses bytes.

  \begin{itemize}
    \item Quem manda primeiro?
    \item Como se pede um documento?
    \item Como se responde que o documento não existe?
  \end{itemize}

  \vspace{0.3cm}
  Um \textbf{protocolo} é esse acordo.

  \vspace{0.3cm}
  Na web, o protocolo é o \textbf{HTTP} (\textit{HyperText Transfer Protocol}),
  criado por Tim Berners-Lee no início dos anos 1990.
\end{frame}

\begin{frame}{Onde o HTTP fica}
  \begin{center}
    \begin{tikzpicture}[
        layer/.style={draw, minimum width=6.5cm, minimum height=0.95cm,
                      align=center, font=\small}
      ]
      \node[layer, fill=blue!12]   (app) at (0,0)    {\textbf{Aplicação}: HTTP, DNS, SMTP};
      \node[layer, fill=gray!10]   (tra) at (0,-1.1) {\textbf{Transporte}: TCP, UDP};
      \node[layer, fill=gray!10]   (red) at (0,-2.2) {\textbf{Rede}: IP};
      \node[layer, fill=gray!10]   (enl) at (0,-3.3) {\textbf{Enlace}: Ethernet, Wi-Fi};
    \end{tikzpicture}
  \end{center}

  O HTTP \textbf{não} se preocupa em fazer os bytes chegarem em ordem e sem
  perdas. Isso é tarefa do TCP, abaixo dele. O HTTP assume um canal confiável e
  define o que trafega por ele.
\end{frame}

\begin{frame}{RFCs}
  As especificações da Internet são publicadas como \textbf{RFCs}
  (\textit{Request for Comments}), numeradas e mantidas pela IETF.

  \begin{alertblock}{Uma RFC nunca é editada depois de publicada}
    Quando a especificação muda, publica-se uma RFC nova que \textbf{torna
    obsoleta} (\textit{obsoletes}) a anterior.
  \end{alertblock}

  O HTTP/1.1 foi definido pela RFC 2616, de 1999. Ela está \tcr{obsoleta desde
  2014}, e a maior parte do material antigo na Internet ainda a cita.
\end{frame}

\begin{frame}{Especificações vigentes}
  \begin{center}
    \small
    \begin{tabular}{@{}llp{4.9cm}@{}}
      \toprule
      \textbf{RFC} & \textbf{Ano} & \textbf{Define} \\
      \midrule
      9110 & 2022 & Semântica do HTTP (métodos, status, cabeçalhos), comum a todas as versões \\
      9111 & 2022 & Cache \\
      9112 & 2022 & Sintaxe do HTTP/1.1 \\
      9113 & 2022 & HTTP/2 \\
      9114 & 2022 & HTTP/3 \\
      \bottomrule
    \end{tabular}
  \end{center}

  \vspace{0.4cm}
  Ao abrir uma RFC, procure no topo os campos \texttt{Obsoletes} e
  \texttt{Obsoleted by}. É assim que se sabe se o documento ainda vale.
\end{frame}

\begin{frame}{Duas formas de ver uma mensagem real}
  Nos próximos minutos você vai olhar mensagens HTTP \textbf{reais}, de dois
  jeitos diferentes:

  \begin{enumerate}
    \item pelo terminal, com o \texttt{curl};
    \item pelo navegador, com as ferramentas de desenvolvimento.
  \end{enumerate}

  \vspace{0.5cm}
  Os detalhes não precisam ser entendidos ainda. Cada um deles é uma seção desta
  aula, e as duas ferramentas voltam em detalhe na parte prática.
\end{frame}

\begin{frame}[fragile]{No terminal}
  \begin{minted}{bash}
curl -v https://example.com
  \end{minted}

  \begin{itemize}
    \item \texttt{>} são as linhas que o \texttt{curl} \textbf{enviou};
    \item \texttt{<} são as linhas que o servidor \textbf{devolveu};
    \item \texttt{*} são comentários do próprio \texttt{curl}.
  \end{itemize}

  \vspace{0.3cm}
  Quatro coisas a notar: a conversa é \textbf{texto}; tem uma primeira linha
  diferente das outras; depois vem uma pilha de pares \texttt{Nome: valor}; por
  fim uma linha em branco, e só então o HTML.
\end{frame}

\begin{frame}{No navegador}
  \begin{enumerate}
    \item Abra qualquer site;
    \item \texttt{F12} ou \texttt{Ctrl+Shift+I};
    \item aba \textbf{Rede} (\textit{Network});
    \item recarregue a página com \texttt{Ctrl+R};
    \item clique em qualquer linha da lista que aparecer.
  \end{enumerate}

  \vspace{0.4cm}
  Na aba \textbf{Cabeçalhos} (\textit{Headers}) você encontra os mesmos dois
  blocos: o que o navegador mandou e o que o servidor respondeu.

  \vspace{0.3cm}
  \begin{block}{}
    É a mesma conversa que o \texttt{curl} mostrou. Muda a apresentação.
  \end{block}
\end{frame}

\begin{frame}{A mesma mensagem nas duas ferramentas}
  \begin{center}
    \small
    \begin{tabular}{@{}lll@{}}
      \toprule
      & \textbf{No \texttt{curl -v}} & \textbf{No navegador} \\
      \midrule
      Endereço pedido    & 1\textsuperscript{a} linha após \texttt{>} & \textit{Request URL} \\
      Verbo              & \texttt{GET}, no início    & \textit{Request Method} \\
      Deu certo?         & 1\textsuperscript{a} linha após \texttt{<} & \textit{Status Code} \\
      Metadados enviados & pares após \texttt{>}      & \textit{Request Headers} \\
      Metadados recebidos& pares após \texttt{<}      & \textit{Response Headers} \\
      Conteúdo           & depois da linha em branco  & aba \textit{Resposta} \\
      \bottomrule
    \end{tabular}
  \end{center}

  \vspace{0.4cm}
  As seis linhas da tabela são, nesta ordem, o roteiro das próximas seções.
\end{frame}

% ==========================================
\section{A mensagem}
% ==========================================

\begin{frame}[fragile]{Anatomia de uma mensagem}
  Toda mensagem HTTP/1.1, de requisição ou de resposta, tem a mesma forma:

  \begin{minted}[fontsize=\small]{text}
<linha inicial>
<Nome-do-cabecalho>: <valor>
<Nome-do-cabecalho>: <valor>
<linha em branco>
<corpo, opcional>
  \end{minted}

  \vspace{0.3cm}
  E o mais importante: \textbf{é texto}. Você pode digitá-la à mão.
\end{frame}

\begin{frame}[fragile]{CRLF e a linha em branco}
  \begin{block}{CRLF}
    Cada linha termina com \textbf{dois} caracteres: \verb|\r\n| (retorno de
    carro e nova linha). Não é apenas \verb|\n|, como é usual em arquivos de
    texto no Linux. Terminador errado, mensagem rejeitada.
  \end{block}

  \begin{block}{A linha em branco}
    Separa os cabeçalhos do corpo. É \textbf{obrigatória}, mesmo quando não há
    corpo.

    \vspace{0.2cm}
    É por isso que, ao digitar uma requisição à mão, você pressiona
    \texttt{Enter} duas vezes.
  \end{block}
\end{frame}

\begin{frame}[fragile]{Requisição}
  \begin{minted}{http}
GET /docs/index.html HTTP/1.1
Host: www.example.com
User-Agent: curl/8.5.0
Accept: text/html

  \end{minted}

  A \textbf{linha inicial} tem três partes separadas por espaço:
  \begin{enumerate}
    \item o \tcb{método} (\texttt{GET}): o que se quer fazer;
    \item o \tcb{alvo} (\texttt{/docs/index.html}): o caminho do recurso;
    \item a \tcb{versão} do protocolo (\texttt{HTTP/1.1}).
  \end{enumerate}

  \vspace{0.2cm}
  Repare: o alvo \textbf{não contém o nome do site}.
\end{frame}

\begin{frame}[fragile]{Resposta}
  \begin{minted}{http}
HTTP/1.1 200 OK
Date: Mon, 23 May 2026 22:38:34 GMT
Server: nginx
Content-Type: text/html; charset=UTF-8
Content-Length: 47

<html>
<body>
<p>Ola mundo</p>
</body>
</html>
  \end{minted}

  \begin{enumerate}
    \item a \tcb{versão} do protocolo;
    \item o \tcb{código de status} (\texttt{200}), para o programa;
    \item a \tcb{frase de motivo} (\texttt{OK}), para o humano, ignorada pelo programa.
  \end{enumerate}
\end{frame}

\begin{frame}{O cabeçalho Host}
  \texttt{Host} é obrigatório em HTTP/1.1 e não existia em HTTP/1.0.
  A razão é econômica.

  \vspace{0.4cm}
  Um único servidor, com um único endereço IP, pode hospedar centenas de sites. Quando
  a requisição chega, o servidor precisa decidir \textbf{qual site entregar}, e a
  única informação disponível é o cabeçalho \texttt{Host}.

  \vspace{0.4cm}
  Chama-se \textbf{hospedagem virtual por nome}
  (\textit{name-based virtual hosting}).
\end{frame}

\begin{frame}{DNS, TCP e TLS antes do HTTP}
  \begin{enumerate}
    \item O navegador consulta o \textbf{DNS} para descobrir o IP de
          \texttt{www.example.com};
    \item Abre uma \textbf{conexão TCP} com esse IP;
    \item (se HTTPS) executa o \textbf{handshake TLS};
    \item Envia a requisição, informando \textbf{dentro dela} com qual nome
          pretendia falar.
  \end{enumerate}

  \vspace{0.5cm}
  O passo 4 é o que torna a hospedagem virtual possível.
\end{frame}

\begin{frame}{Onde a mensagem termina}
  Pergunta que parece trivial e não é:

  \begin{block}{}
    Recebida a última linha de cabeçalho, como o cliente sabe \textbf{em que
    ponto} a resposta acabou?
  \end{block}

  \vspace{0.4cm}
  O TCP entrega um fluxo contínuo de bytes. Ele \textbf{não} sinaliza fronteiras
  de mensagem.

  \vspace{0.4cm}
  O HTTP resolve isso de três formas.
\end{frame}

\begin{frame}[fragile]{1. Content-Length}
  O servidor informa o tamanho exato do corpo em bytes. O cliente conta.

  \begin{minted}{http}
HTTP/1.1 200 OK
Content-Length: 47

  \end{minted}

  \vspace{0.2cm}
  Exige que o servidor conheça o tamanho total \textbf{antes} de começar a
  enviar.
\end{frame}

\begin{frame}[fragile]{2. Transfer-Encoding: chunked}
  O corpo vai em blocos, cada um precedido pelo seu tamanho em hexadecimal.
  Um bloco de tamanho \texttt{0} encerra a mensagem.

  \begin{minted}{text}
HTTP/1.1 200 OK
Transfer-Encoding: chunked

7
Ola mun
3
do!
0

  \end{minted}

  Corpo reconstruído: \texttt{Ola mundo}.

  \vspace{0.2cm}
  Permite responder \textbf{antes} de saber o tamanho final. É a base do envio
  contínuo de dados (\textit{streaming}).
\end{frame}

\begin{frame}{3. Fechamento da conexão}
  O servidor envia \texttt{Connection: close}, transmite o corpo e encerra a
  conexão TCP. O fim da conexão é o fim da mensagem.

  \vspace{0.4cm}
  Era o mecanismo do HTTP/1.0. É ineficiente, porque impede reaproveitar a
  conexão.

  \vspace{0.6cm}
  \begin{alertblock}{Requisição pendurada e resposta truncada}
    Requisição que fica pendurada indefinidamente: o cliente ainda espera bytes
    prometidos por um \texttt{Content-Length} maior que o corpo real.

    \vspace{0.2cm}
    Resposta truncada: o contrário.
  \end{alertblock}
\end{frame}

% ==========================================
% \subsection{URL}
% ==========================================

\begin{frame}[fragile]{Anatomia de uma URL}
  \begin{center}
    \begin{minted}[fontsize=\tiny, breaklines=false, frame=none]{text}
https://usuario@www.example.com:8080/produtos/lista.php?cat=livros&pag=2#topo
\___/   \_____/ \_____________/ \__/\_________________/ \______________/ \__/
  |        |           |          |          |                  |          |
esquema  userinfo    host       porta     caminho           consulta   fragmento
    \end{minted}
  \end{center}

  \vspace{0.5cm}
  Sete partes. Duas delas, \textbf{consulta} e \textbf{fragmento}, são as que
  mais confundem, e são justamente as que interessam nesta aula.
\end{frame}

\begin{frame}{As partes da URL}
  \begin{center}
    \footnotesize
    \begin{tabular}{@{}lp{7.2cm}@{}}
      \toprule
      \textbf{esquema}  & O protocolo: \texttt{http}, \texttt{https}, \texttt{mailto}, \texttt{file}. Determina tudo o que vem depois. \\
      \textbf{userinfo} & Credencial embutida. Em desuso no HTTP e desencorajado por segurança. \\
      \textbf{host}     & Nome de domínio ou IP do servidor. \\
      \textbf{porta}    & Porta TCP. Omitida, vale 80 para \texttt{http} e 443 para \texttt{https}. \\
      \textbf{caminho}  & Identifica o recurso dentro do servidor. \\
      \textbf{consulta} & Pares \texttt{chave=valor} separados por \texttt{\&}, iniciados por \texttt{?}. É o que o PHP lerá em \texttt{\$\_GET}. \\
      \textbf{fragmento}& Iniciado por \texttt{\#}. Identifica parte interna do documento. \\
      \bottomrule
    \end{tabular}
  \end{center}
\end{frame}

\begin{frame}{O fragmento nunca é enviado ao servidor}
  \begin{center}
    \Large \texttt{https://site.com/pagina\#secao}
  \end{center}

  \vspace{0.6cm}
  O trecho \texttt{\#secao} é processado \textbf{inteiramente pelo navegador},
  que o usa para rolar a página até o elemento correspondente.

  \vspace{0.4cm}
  O servidor \textbf{jamais} toma conhecimento dele.

  \vspace{0.6cm}
  \begin{block}{}
    Você vai comprovar isso experimentalmente na parte prática desta aula.
  \end{block}
\end{frame}

\begin{frame}{Codificação percentual}
  Espaços, acentos e os próprios caracteres \texttt{?}, \texttt{\&} e
  \texttt{\#} não podem aparecer cruamente numa URL.

  \vspace{0.3cm}
  São substituídos por \texttt{\%} seguido do valor hexadecimal de cada byte na
  codificação UTF-8.

  \begin{center}
    \small
    \begin{tabular}{@{}ll@{}}
      \toprule
      \textbf{Caractere} & \textbf{Codificado} \\
      \midrule
      espaço  & \texttt{\%20} \\
      \texttt{á} & \texttt{\%C3\%A1} \\
      \texttt{\&} & \texttt{\%26} \\
      \texttt{\#} & \texttt{\%23} \\
      \bottomrule
    \end{tabular}
  \end{center}

  \vspace{0.2cm}
  A busca por \texttt{análise e projeto} vira:

  \begin{center}
    \small\texttt{q=an\%C3\%A1lise\%20e\%20projeto}
  \end{center}

  \vspace{0.2cm}
  \footnotesize Note que \texttt{á} ocupa \textbf{dois} bytes em UTF-8.
\end{frame}

% ==========================================
\section{Semântica}
% ==========================================

\begin{frame}{Passagem de parâmetros}
  Até agora o cliente só pediu: ``me entregue o recurso deste caminho''.

  \vspace{0.4cm}
  Mas uma aplicação web precisa que o cliente \textbf{envie dados}:

  \begin{itemize}
    \item o termo digitado na busca;
    \item a página da listagem que se quer ver;
    \item o login e a senha;
    \item os campos de um cadastro.
  \end{itemize}

  \vspace{0.4cm}
  \begin{block}{}
    Sem isso não existe aplicação, só um servidor de arquivos. Toda a
    disciplina, daqui em diante, gira em torno disso.
  \end{block}
\end{frame}

\begin{frame}{Caminho, consulta e corpo}
  A mensagem HTTP tem três lugares onde dados cabem.

  \begin{center}
    \small
    \begin{tabular}{@{}lll@{}}
      \toprule
      \textbf{Lugar} & \textbf{Exemplo} & \textbf{Onde fica na mensagem} \\
      \midrule
      Caminho   & \texttt{/produtos/42}       & linha inicial \\
      Consulta  & \texttt{/busca?q=teclado}   & linha inicial, após \texttt{?} \\
      Corpo     & \texttt{q=teclado}          & depois da linha em branco \\
      \bottomrule
    \end{tabular}
  \end{center}

  \vspace{0.4cm}
  Cabeçalhos também carregam dados (\texttt{Cookie}, \texttt{Authorization}),
  mas são metadados, e não parâmetros da operação.
\end{frame}

\begin{frame}[fragile]{Parâmetros na própria URL}
  \begin{center}
    \large \texttt{/busca?q=teclado\&pag=2\&ordem=preco}
  \end{center}

  \vspace{0.4cm}
  Regras da \textit{query string}:

  \begin{itemize}
    \item começa no primeiro \texttt{?};
    \item cada parâmetro é um par \texttt{chave=valor};
    \item pares são separados por \texttt{\&};
    \item caracteres especiais vão codificados: \texttt{q=ar\%20condicionado}.
  \end{itemize}

  \vspace{0.3cm}
  O servidor recebe isto e monta uma estrutura de chave e valor. Em PHP:

  \begin{minted}[fontsize=\small]{php}
$_GET['q']      // "teclado"
$_GET['pag']    // "2"
  \end{minted}
\end{frame}

\begin{frame}[fragile]{Parâmetros no corpo da mensagem}
  \begin{minted}{http}
POST /busca HTTP/1.1
Host: loja.com
Content-Type: application/x-www-form-urlencoded
Content-Length: 15

q=teclado&pag=2
  \end{minted}

  \vspace{0.3cm}
  \textbf{Exatamente a mesma gramática} da consulta, só que depois da linha em
  branco. O nome do formato diz isso: \texttt{x-www-form-urlencoded}, ou seja,
  ``codificado como numa URL''.

  \vspace{0.2cm}
  Em PHP, muda apenas o nome da variável: \texttt{\$\_POST['q']}.
\end{frame}

\begin{frame}[fragile]{O formulário HTML}
  Você raramente monta a URL à mão. Quem monta é o formulário HTML.

  \begin{minted}[fontsize=\small]{html}
<form action="/busca" method="get">
  <input name="q">
  <input name="pag">
</form>
  \end{minted}

  \vspace{0.2cm}
  Com \texttt{method="get"}, o navegador coloca os campos na \textbf{consulta}:
  \texttt{/busca?q=teclado\&pag=2}

  \vspace{0.2cm}
  Com \texttt{method="post"}, coloca os mesmos campos no \textbf{corpo}.

  \vspace{0.4cm}
  \begin{alertblock}{GET ou POST?}
    Se os dois transportam os mesmos dados, qual escolher, e por quê?
  \end{alertblock}
\end{frame}

\begin{frame}{Métodos}
  O método é o verbo da requisição.

  \begin{center}
    \small
    \begin{tabular}{@{}lp{7cm}@{}}
      \toprule
      \texttt{GET}     & Solicita a representação de um recurso. \\
      \texttt{POST}    & Envia dados para serem processados pelo recurso. \\
      \texttt{PUT}     & Substitui integralmente o recurso. \\
      \texttt{DELETE}  & Remove o recurso. \\
      \texttt{HEAD}    & Igual ao \texttt{GET}, mas devolve só os cabeçalhos. \\
      \texttt{OPTIONS} & Pergunta quais métodos o servidor aceita. \\
      \texttt{PATCH}   & Altera parcialmente o recurso. \\
      \bottomrule
    \end{tabular}
  \end{center}
\end{frame}

\begin{frame}{GET e POST não se distinguem pelo lugar dos dados}
  \begin{center}
    \large ``GET manda os dados pela URL e POST manda pelo corpo.''
  \end{center}

  \vspace{0.4cm}
  A formulação é verdadeira, mas descreve o \textbf{efeito}, não a
  \textbf{causa}, e pode levar a decisões erradas de projeto.

  \vspace{0.4cm}
  A distinção que importa está em duas propriedades definidas na RFC 9110.
\end{frame}

\begin{frame}{Seguro e idempotente}
  \begin{block}{Método seguro (\textit{safe})}
    Não altera o estado do servidor. É apenas leitura.
  \end{block}

  \begin{block}{Método idempotente}
    Executá-lo uma vez ou N vezes seguidas produz o \textbf{mesmo estado final}
    no servidor.
  \end{block}

  \begin{center}
    \small
    \begin{tabular}{@{}lcc@{}}
      \toprule
      \textbf{Método} & \textbf{Seguro} & \textbf{Idempotente} \\
      \midrule
      \texttt{GET}    & \tcg{sim} & \tcg{sim} \\
      \texttt{HEAD}   & \tcg{sim} & \tcg{sim} \\
      \texttt{PUT}    & \tcr{não} & \tcg{sim} \\
      \texttt{DELETE} & \tcr{não} & \tcg{sim} \\
      \texttt{POST}   & \tcr{não} & \tcr{não} \\
      \bottomrule
    \end{tabular}
  \end{center}

  \vspace{0.2cm}
  \footnotesize
  \texttt{DELETE} não é seguro, mas é idempotente: apagar o mesmo recurso duas
  vezes deixa o servidor no mesmo estado, ainda que a segunda resposta seja
  \texttt{404}. \texttt{POST} não é, porque enviar duas vezes o formulário de
  compra cria duas compras.
\end{frame}

\begin{frame}{Consequências no navegador}
  \begin{itemize}
    \item Ele pode buscar antecipadamente um link, porque \texttt{GET} é seguro
          e não causa dano;
    \item Ele recarrega um \texttt{GET} sem perguntar nada, e exige confirmação
          antes de reenviar um \texttt{POST};
    \item Caches e servidores intermediários armazenam respostas de
          \texttt{GET}, e nunca de \texttt{POST};
    \item O histórico guarda a URL, portanto guarda os dados de um \texttt{GET}.
  \end{itemize}
\end{frame}

\begin{frame}{Senhas e links que apagam}
  \begin{alertblock}{Senha não viaja por GET}
    Mesmo em HTTPS. Ela ficaria no histórico do navegador e nos registros de
    acesso do servidor.
  \end{alertblock}

  \begin{alertblock}{Link que apaga registro não é GET}
    Qualquer mecanismo que percorra links da página o executaria sem intenção.
  \end{alertblock}
\end{frame}

\begin{frame}{Códigos de status: as classes}
  O primeiro dígito define a classe.

  \begin{center}
    \small
    \begin{tabular}{@{}lp{7cm}@{}}
      \toprule
      \texttt{1xx} & Informativo. Processamento em andamento. \\
      \texttt{2xx} & Sucesso. \\
      \texttt{3xx} & Redirecionamento. O cliente precisa de uma ação adicional. \\
      \texttt{4xx} & Erro do \textbf{cliente}. A requisição está errada. \\
      \texttt{5xx} & Erro do \textbf{servidor}. A requisição estava certa. \\
      \bottomrule
    \end{tabular}
  \end{center}

  \vspace{0.4cm}
  \begin{block}{A fronteira entre 4xx e 5xx indica onde procurar o erro}
    Ao depurar um sistema, \texttt{4xx} manda olhar o código que fez a
    requisição, e \texttt{5xx} manda olhar o servidor.
  \end{block}
\end{frame}

\begin{frame}{Códigos frequentes}
  \begin{center}
    \scriptsize
    \begin{tabular}{@{}llp{5.6cm}@{}}
      \toprule
      \textbf{Código} & \textbf{Nome} & \textbf{Quando ocorre} \\
      \midrule
      \texttt{200} & OK                    & Sucesso, com corpo. \\
      \texttt{201} & Created               & Recurso criado. Acompanha \texttt{Location}. \\
      \texttt{204} & No Content            & Sucesso sem corpo. Não pode ter corpo. \\
      \texttt{301} & Moved Permanently     & Mudou definitivamente. O navegador memoriza. \\
      \texttt{302} & Found                 & Mudou temporariamente. Não memoriza. \\
      \texttt{304} & Not Modified          & Não mudou desde a última vez. Sem corpo. \\
      \texttt{400} & Bad Request           & Requisição malformada. \\
      \texttt{401} & Unauthorized          & Falta autenticação. Acompanha \texttt{WWW-Authenticate}. \\
      \texttt{403} & Forbidden             & Entendeu, sabe quem você é, e recusa. \\
      \texttt{404} & Not Found             & O recurso não existe naquele caminho. \\
      \texttt{405} & Method Not Allowed    & Recurso existe, método não é aceito. Acompanha \texttt{Allow}. \\
      \texttt{500} & Internal Server Error & Falha no servidor. Em PHP, erro fatal no script. \\
      \bottomrule
    \end{tabular}
  \end{center}
\end{frame}

\begin{frame}{401 contra 403, 301 contra 302}
  \begin{block}{401 contra 403}
    \texttt{401}: ``identifique-se''. \\
    \texttt{403}: ``já sei quem você é e ainda assim não pode''.
  \end{block}

  \begin{block}{301 contra 302}
    O \texttt{301} é memorizado pelo navegador de forma persistente.

    \vspace{0.2cm}
    Publicar um \texttt{301} errado é \textbf{difícil de reverter}: os
    navegadores dos usuários continuarão indo ao endereço antigo mesmo depois
    da correção.
  \end{block}
\end{frame}

% ==========================================
% \subsection{Cabeçalhos}
% ==========================================

\begin{frame}{Cabeçalhos}
  Pares \texttt{Nome: valor} que carregam os metadados da mensagem.
  O nome não diferencia maiúsculas de minúsculas. Existem centenas.

  \begin{columns}[t]
    \begin{column}{0.48\textwidth}
      \textbf{Enviados pelo cliente}
      \scriptsize
      \begin{description}
        \item[Host] Nome do site. Obrigatório.
        \item[User-Agent] Identifica o programa.
        \item[Accept] Tipos que sabe processar.
        \item[Accept-Language] Idiomas preferidos.
        \item[Accept-Encoding] Compressões.
        \item[Content-Type] Formato do corpo.
        \item[Content-Length] Tamanho do corpo.
        \item[Cookie] Dados devolvidos.
        \item[Authorization] Credenciais.
      \end{description}
    \end{column}
    \begin{column}{0.48\textwidth}
      \textbf{Enviados pelo servidor}
      \scriptsize
      \begin{description}
        \item[Content-Type] Formato do corpo.
        \item[Content-Length] Tamanho do corpo.
        \item[Location] Destino em \texttt{3xx}.
        \item[Set-Cookie] Pede para armazenar.
        \item[Cache-Control] Regras de cache.
        \item[ETag] Versão atual do recurso.
        \item[Server] Software servidor.
      \end{description}
    \end{column}
  \end{columns}
\end{frame}

\begin{frame}{Negociação de conteúdo}
  Os cabeçalhos \texttt{Accept*} permitem que cliente e servidor combinem o
  formato da resposta.

  \vspace{0.3cm}
  O cliente \textbf{declara} o que aceita. O servidor \textbf{escolhe} entre o
  que sabe produzir.

  \vspace{0.4cm}
  A mesma URL pode devolver HTML para um navegador e JSON para um programa.

  \vspace{0.5cm}
  \begin{alertblock}{Nem todo servidor negocia}
    Negociação de conteúdo é um recurso do protocolo. Cabe a cada servidor
    implementá-la ou não. Muitos escolhem o formato pelo caminho da URL.
  \end{alertblock}
\end{frame}

\begin{frame}{Content-Type}
  Informa como interpretar os bytes do corpo. Segue o padrão MIME.

  \begin{center}
    \footnotesize
    \begin{tabular}{@{}ll@{}}
      \toprule
      \texttt{text/html}    & Documento HTML \\
      \texttt{text/css}     & Folha de estilo \\
      \texttt{application/javascript} & Script \\
      \texttt{application/json}       & Dados em JSON \\
      \texttt{image/png}              & Imagem PNG \\
      \texttt{application/x-www-form-urlencoded} & Formulário HTML padrão \\
      \texttt{multipart/form-data}    & Formulário com upload de arquivo \\
      \bottomrule
    \end{tabular}
  \end{center}
\end{frame}

\begin{frame}{charset e os acentos quebrados}
  \begin{center}
    \large \texttt{Content-Type: text/html; charset=UTF-8}
  \end{center}

  \vspace{0.5cm}
  O parâmetro \texttt{charset} diz que os bytes devem ser lidos como UTF-8.

  \vspace{0.5cm}
  \begin{alertblock}{A causa mais comum de \texttt{ção} na tela}
    Um \texttt{charset} errado ou ausente.

    \vspace{0.2cm}
    O problema quase nunca está no arquivo. Está na \textbf{declaração}.
  \end{alertblock}
\end{frame}

% ==========================================
\section{Além de uma requisição}
% ==========================================

\begin{frame}{Protocolo sem estado}
  \begin{block}{\textit{Stateless}}
    Cada requisição é interpretada de forma isolada. O servidor não guarda
    nenhuma lembrança das requisições anteriores.
  \end{block}

  \vspace{0.4cm}
  Duas requisições do mesmo usuário, com um segundo de intervalo, são para o
  servidor \textbf{dois eventos sem relação entre si}.
\end{frame}

\begin{frame}{Escala, e o problema do usuário logado}
  \begin{itemize}
    \item Sem estado, \textbf{qualquer} servidor em cluster (ou ``fazenda de
          servidores'') pode atender \textbf{qualquer} requisição;
    \item Derrubar um servidor não derruba as sessões em curso.
  \end{itemize}

  Foi o que permitiu a web escalar. Mas:

  \vspace{0.4cm}
  \begin{center}
    \Large Se o servidor não se lembra de nada, \\ como existe ``usuário logado''?
  \end{center}

  \vspace{0.4cm}
  A única saída é o \textbf{cliente reenviar}, em toda requisição, alguma
  informação que permita ao servidor reconstruir o contexto.
\end{frame}

\begin{frame}{Cookies, sessões e tokens}
  \begin{description}
    \item[Cookies] O servidor responde com \texttt{Set-Cookie: nome=valor}. O
      navegador armazena e passa a incluir \texttt{Cookie: nome=valor} em toda
      requisição para aquele domínio. É apenas texto em cabeçalho.

    \item[Sessões] O cookie guarda só um identificador aleatório. Os dados ficam
      no servidor, associados a esse identificador. É o modelo que usaremos com
      PHP.

    \item[Tokens] Uma credencial assinada, enviada em \texttt{Authorization}.
      Comum em APIs.
  \end{description}

  \vspace{0.4cm}
  Todos resolvem o mesmo problema pela mesma via: \textbf{fazer o cliente
  carregar o contexto}, porque o protocolo não o carrega.
\end{frame}

% ==========================================
% \subsection{Versões e conexões}
% ==========================================

\begin{frame}{HTTP/1.0 e HTTP/1.1}
  \begin{block}{HTTP/1.0}
    Uma conexão TCP \textbf{por recurso}. Uma página com trinta imagens exigia
    trinta conexões, cada uma com o custo de abertura e fechamento.
  \end{block}

  \begin{block}{HTTP/1.1}
    \textbf{Conexões persistentes}: a conexão TCP permanece aberta e várias
    requisições são feitas em sequência por ela.
  \end{block}
\end{frame}

\begin{frame}{Bloqueio de cabeça de fila}
  No HTTP/1.1 as requisições são atendidas \textbf{em ordem}. Uma resposta lenta
  bloqueia todas as seguintes na mesma conexão.

  \vspace{0.4cm}
  Chama-se \textbf{bloqueio de cabeça de fila}
  (\textit{head-of-line blocking}).

  \vspace{0.4cm}
  Os navegadores contornavam abrindo cerca de \textbf{seis conexões paralelas}
  por domínio.
\end{frame}

\begin{frame}{HTTP/2 (2015) e HTTP/3 (2022)}
  \begin{block}{HTTP/2}
    Mensagens deixam de ser texto e passam a ser \textbf{binárias}, em quadros.
    Traz \textbf{multiplexação}: várias requisições e respostas trafegam
    simultaneamente e intercaladas na mesma conexão, eliminando o bloqueio no
    nível do HTTP. Comprime cabeçalhos com HPACK.
  \end{block}

  \begin{block}{HTTP/3}
    Substitui o TCP pelo \textbf{QUIC}, que roda sobre UDP e implementa
    confiabilidade por conta própria. Elimina o bloqueio de cabeça de fila que
    \textbf{persistia na camada TCP} e funde o estabelecimento de conexão com o
    handshake criptográfico, reduzindo as idas e vindas até o primeiro byte.
  \end{block}
\end{frame}

\begin{frame}{A semântica é a mesma nas três versões}
  Métodos, códigos de status e cabeçalhos são definidos na RFC 9110 e valem para
  HTTP/1.1, HTTP/2 e HTTP/3.

  \vspace{0.4cm}
  O que muda é a \textbf{codificação} e o \textbf{transporte}.

  \vspace{0.6cm}
  \begin{block}{HTTP/1.1 é o único legível por humanos}
    Porque ele é texto e pode ser digitado à mão.

    \vspace{0.2cm}
    A versão 2 é a mesma conversa, em formato ilegível para humanos.
  \end{block}
\end{frame}

% ==========================================
% \subsection{HTTPS e origem}
% ==========================================

\begin{frame}{HTTPS}
  HTTPS é HTTP transportado dentro de um túnel \textbf{TLS}
  (\textit{Transport Layer Security}, sucessor do SSL).

  \vspace{0.4cm}
  O protocolo HTTP \textbf{não muda}. As mesmas mensagens de texto trafegam,
  agora cifradas.

  \vspace{0.5cm}
  No \textit{handshake}, cliente e servidor negociam algoritmos, o servidor
  apresenta seu \textbf{certificado digital}, e ambos derivam uma chave de
  sessão usada para cifrar o restante da conversa.
\end{frame}

\begin{frame}{O que o TLS garante}
  \begin{description}
    \item[Confidencialidade] Quem observa a rede não lê o conteúdo das
      mensagens: caminho da URL, cabeçalhos, cookies e corpo.
    \item[Integridade] Alterações no tráfego são detectadas.
    \item[Autenticação do servidor] O certificado, assinado por uma autoridade
      certificadora reconhecida pelo navegador, comprova que o servidor controla
      aquele nome de domínio.
  \end{description}
\end{frame}

\begin{frame}{O que o TLS não garante}
  \begin{itemize}
    \item \textbf{Não esconde com quem você fala.} O IP de destino é visível, e
          o nome do host viaja em claro no handshake (campo SNI);
    \item \textbf{Não atesta idoneidade.} Certificados são gratuitos e obtidos em
          minutos. Um site fraudulento tem cadeado igual ao de um banco;
    \item \textbf{Não protege o dado depois que ele chega.} Segurança do banco de
          dados, das senhas e do código do servidor são problemas separados.
  \end{itemize}

  \vspace{0.3cm}
  \begin{alertblock}{Conteúdo misto}
    Página HTTPS que carrega script por HTTP: o atacante troca o script e
    controla a página inteira. A garantia é anulada pelo elo mais fraco, e por
    isso os navegadores bloqueiam a carga.
  \end{alertblock}
\end{frame}

\begin{frame}{Origem}
  \begin{center}
    \Large \textbf{esquema + host + porta}
  \end{center}

  \vspace{0.4cm}
  Duas URLs são da mesma origem somente se as \textbf{três} coincidirem.

  \begin{center}
    \scriptsize
    \begin{tabular}{@{}llc@{}}
      \toprule
      \textbf{URL A} & \textbf{URL B} & \textbf{Mesma origem?} \\
      \midrule
      \texttt{https://site.com/a}  & \texttt{https://site.com/b/c}    & \tcg{sim}, caminho não conta \\
      \texttt{http://site.com}     & \texttt{https://site.com}        & \tcr{não}, esquema difere \\
      \texttt{https://site.com}    & \texttt{https://api.site.com}    & \tcr{não}, host difere \\
      \texttt{http://site.com}     & \texttt{http://site.com:8080}    & \tcr{não}, porta difere \\
      \bottomrule
    \end{tabular}
  \end{center}
\end{frame}

\begin{frame}{Política de mesma origem e CORS}
  Regra imposta pelo \textbf{navegador}: JavaScript executando numa origem não
  pode ler a resposta de uma requisição feita a outra origem. Sem ela, uma
  página maliciosa faria requisições ao seu banco em outra aba, aproveitando os
  cookies já armazenados, e leria o resultado.

  \vspace{0.3cm}
  \textbf{CORS} (\textit{Cross-Origin Resource Sharing}) é como o servidor de
  destino autoriza essa leitura, pelo cabeçalho
  \texttt{Access-Control-Allow-Origin}. Para requisições que possam alterar
  dados, o navegador antes envia uma verificação prévia com \texttt{OPTIONS}.

  \begin{block}{Quem impõe e quem autoriza}
    \begin{enumerate}
      \item A restrição é do \textbf{navegador}, não do protocolo. O
            \texttt{curl} acessa qualquer origem. Por isso uma URL funciona no
            terminal e falha no console.
      \item Quem autoriza é o \textbf{servidor de destino}. Não há nada no seu
            JavaScript que contorne a recusa.
    \end{enumerate}
  \end{block}
\end{frame}

% ==========================================
% \subsection{Cache}
% ==========================================

\begin{frame}{Cache por prazo}
  \begin{center}
    \large \texttt{Cache-Control: max-age=3600}
  \end{center}

  \vspace{0.5cm}
  O servidor declara que a resposta vale por 3600 segundos.

  \vspace{0.4cm}
  Durante esse período o navegador usa a cópia local e \textbf{não faz
  requisição nenhuma}.
\end{frame}

\begin{frame}{Requisição condicional}
  Vencido o prazo, o cliente não precisa baixar tudo de novo. Ele
  \textbf{pergunta se mudou}, usando um identificador de versão que guardou:

  \begin{description}
    \item[ETag] Etiqueta opaca que identifica a versão do recurso. \\
      Servidor: \texttt{ETag: "abc123"} \\
      Cliente reenvia: \texttt{If-None-Match: "abc123"}
    \item[Last-Modified] Data da última alteração. \\
      Cliente reenvia: \texttt{If-Modified-Since: <data>}
  \end{description}

  \vspace{0.4cm}
  Se não mudou, o servidor responde \texttt{304 Not Modified},
  \textbf{sem corpo}.

  \vspace{0.2cm}
  A economia é o corpo inteiro. Só os cabeçalhos trafegam.
\end{frame}

% ==========================================
\section{Prática}
% ==========================================

\begin{frame}{Preparando o ambiente}
  A ferramenta principal é o \textbf{\texttt{curl}}, um cliente HTTP de linha de
  comando.

  \begin{itemize}
    \item \textbf{Linux e macOS}: já instalado.
    \item \textbf{Windows 10 e 11}: já instalado, como \texttt{curl.exe}.
  \end{itemize}

  \vspace{0.4cm}
  \begin{alertblock}{Atenção, usuários de Windows}
    No PowerShell, \texttt{curl} é apelido de \texttt{Invoke-WebRequest}, que
    tem sintaxe completamente diferente. Escreva \texttt{curl.exe} com a
    extensão.

    \vspace{0.2cm}
    No \texttt{cmd} e no Git Bash, \texttt{curl} funciona normalmente.
    Recomenda-se o \textbf{Git Bash} ou um ambiente \textbf{WSL}.
  \end{alertblock}
\end{frame}

\begin{frame}[fragile]{Ver a conversa completa}
  \begin{minted}{bash}
curl -v https://example.com
  \end{minted}

  A opção \texttt{-v} (\textit{verbose}) exibe a mensagem trocada:

  \begin{itemize}
    \item \texttt{>} linhas \textbf{enviadas};
    \item \texttt{<} linhas \textbf{recebidas};
    \item \texttt{*} comentários do próprio \texttt{curl} sobre a conexão.
  \end{itemize}

  \vspace{0.3cm}
  Identifique na saída: resolução do nome, handshake TLS, linha de requisição,
  cabeçalhos enviados, linha de status, cabeçalhos recebidos, corpo.
\end{frame}

\begin{frame}[fragile]{Somente os cabeçalhos}
  \begin{minted}{bash}
curl -s -o /dev/null -D - https://developer.mozilla.org
  \end{minted}

  \begin{itemize}
    \item \texttt{-s} silencia a barra de progresso;
    \item \texttt{-o /dev/null} descarta o corpo;
    \item \texttt{-D -} escreve os cabeçalhos na saída padrão.
  \end{itemize}

  \begin{minted}{bash}
curl -I https://developer.mozilla.org
  \end{minted}

  A opção \texttt{-I} faz algo parecido, porém enviando o método
  \texttt{HEAD} em vez de \texttt{GET}. São requisições \textbf{distintas}, e
  alguns servidores respondem diferente a cada uma.
\end{frame}

\begin{frame}[fragile]{Enviar dados: formulário}
  \begin{minted}{bash}
curl -v -d "nome=Ana&curso=TADS" https://postman-echo.com/post
  \end{minted}

  A opção \texttt{-d}, sozinha, define o método como \texttt{POST} e o
  cabeçalho \texttt{Content-Type: application/x-www-form-urlencoded}.

  \vspace{0.4cm}
  \begin{block}{}
    Este é exatamente o formato que o PHP lerá em \texttt{\$\_POST}, na segunda
    metade da disciplina.
  \end{block}
\end{frame}

\begin{frame}[fragile]{Cookies com \texttt{-c} e \texttt{-b}}
  \begin{minted}{bash}
curl -s -o /dev/null -c jar.txt \
     "https://postman-echo.com/cookies/set?turma=ds122"
cat jar.txt
curl -s -b jar.txt https://postman-echo.com/cookies
curl -s           https://postman-echo.com/cookies
  \end{minted}

  \begin{itemize}
    \item \texttt{-c jar.txt} grava os cookies recebidos, fazendo o papel do
          armazenamento do navegador;
    \item \texttt{-b jar.txt} reenvia esses cookies.
  \end{itemize}
\end{frame}

\begin{frame}[fragile]{Com e sem o cookie}
  \begin{minted}[fontsize=\scriptsize]{text}
com -b:   {"cookies":{ ... ,"turma":"ds122"}}
sem -b:   {"cookies":{}}
  \end{minted}

  \vspace{0.4cm}
  O último comando é idêntico ao anterior, exceto pela ausência de \texttt{-b}.

  \begin{center}
    \Large O servidor não guardou nada.
  \end{center}

  \vspace{0.3cm}
  \footnotesize
  O primeiro comando responde \texttt{302}, porque o serviço redireciona depois
  de gravar. E o arquivo \texttt{jar.txt} conterá também cookies definidos pela
  infraestrutura do site.
\end{frame}

\begin{frame}[fragile]{Ser o servidor}
  Até aqui você foi o cliente. Agora inverta os papéis com o \texttt{nc}
  (\textit{netcat}), que abre conexões TCP brutas.

  \begin{minted}{bash}
nc -l 8080          # macOS, OpenBSD
nc -l -p 8080       # GNU netcat
ncat -l 8080        # Fedora e derivados (nmap-ncat)
  \end{minted}

  Com o comando rodando e o terminal aparentemente travado, abra no navegador:

  \begin{minted}{text}
http://localhost:8080/pagina?a=1&b=2#secao
  \end{minted}
\end{frame}

\begin{frame}{O que aparece na requisição do navegador}
  \begin{enumerate}
    \item O navegador envia \textbf{muito mais cabeçalhos} que o \texttt{curl}:
          \texttt{User-Agent}, \texttt{Accept}, \texttt{Accept-Encoding},
          \texttt{Accept-Language}, \texttt{Sec-Fetch-*} e outros.
          Compare com a saída do \texttt{curl -v}.

    \item A consulta \texttt{?a=1\&b=2} aparece na linha de requisição.
          O fragmento \texttt{\#secao} \textbf{não aparece em lugar nenhum}.

    \item O navegador continua carregando, porque \textbf{ninguém respondeu}.
  \end{enumerate}
\end{frame}

\begin{frame}[fragile]{Responda à mão}
  Digite no terminal do \texttt{nc} e pressione \texttt{Enter} duas vezes após a
  linha \texttt{Content-Length}:

  \begin{minted}{http}
HTTP/1.1 200 OK
Content-Type: text/html; charset=utf-8
Content-Length: 38

<h1>Resposta digitada byte a byte</h1>
  \end{minted}

  Confira: \texttt{<h1>} tem 4 bytes, o texto tem 29, \texttt{</h1>} tem 5.
  Total 38. Se errar para mais, o navegador ficará aguardando bytes que nunca
  chegam.

  \vspace{0.3cm}
  \begin{center}
    \large Você acabou de operar como servidor web.
  \end{center}
\end{frame}

\begin{frame}{De volta às ferramentas do navegador}
  Abrimos a aba \textbf{Rede} no começo da aula sem saber o nome de nada.
  Volte nela agora: cada coluna e cada campo tem nome, e você já os conhece.

  \vspace{0.3cm}
  Recomendados: \textbf{Firefox} e navegadores baseados em \textbf{Chromium}.

  \vspace{0.5cm}
  \begin{block}{Copiar como cURL}
    Na aba Rede, clique com o botão direito sobre uma requisição e escolha
    \textbf{Copiar como cURL}. Cole no terminal e execute.

    \vspace{0.2cm}
    Mesma resposta. Navegador e \texttt{curl} falam exatamente o mesmo
    protocolo. A diferença está na quantidade de cabeçalhos.
  \end{block}
\end{frame}

% ==========================================
\section{Exercícios}
% ==========================================

\begin{frame}{Para fazer durante a semana}
  \begin{block}{Sem entrega e sem nota}
    Estes exercícios são de \textbf{fixação}. Nada é recolhido nesta aula.
  \end{block}

  \vspace{0.4cm}
  Ainda assim, faça-os. O conteúdo desta aula cai na Prova 1, e digitar os
  comandos é a única forma de fixar o que vimos aqui.

  \vspace{0.4cm}
  Sugestão de método: para cada item, anote o \textbf{comando} usado, o
  \textbf{trecho relevante} da saída e a \textbf{resposta} da pergunta. Guarde
  num arquivo seu. Vai servir de resumo na semana da prova.

  \vspace{0.4cm}
  Dúvidas na aula que vem, ou no fórum da UFPR Virtual.
\end{frame}

\begin{frame}{Exercícios 1 a 3}
  \begin{enumerate}
    \item Escolha três sites. Para cada um, registre o código de status, a
          versão do HTTP e o valor de \texttt{Server}, quando presente. Explique
          as diferenças observadas.

    \item Encontre uma URL que responda \texttt{301} e outra que responda
          \texttt{404}. Para a primeira, mostre a cadeia completa de
          redirecionamento e informe quantos saltos foram necessários.

    \item Envie \texttt{nome=SeuNome\&periodo=2} por \texttt{GET} e por
          \texttt{POST} para um serviço de eco. Mostre onde os dados aparecem em
          cada caso. Explique por que a URL do \texttt{GET} fica no histórico e
          a do \texttt{POST} não, relacionando com segurança e idempotência.
  \end{enumerate}
\end{frame}

\begin{frame}{Exercícios 4 a 7}
  \begin{enumerate}
    \setcounter{enumi}{3}
    \item Demonstre experimentalmente que o HTTP não guarda estado, usando
          \texttt{-c} e \texttt{-b}. Descreva, em no máximo três frases, o que o
          servidor recebeu em cada requisição.

    \item Obtenha o \texttt{ETag} de um recurso, reenvie-o em
          \texttt{If-None-Match} e registre o código e o tamanho do corpo da
          segunda resposta. Explique a economia produzida.

    \item Execute o \texttt{nc} como servidor, acesse
          \texttt{http://localhost:8080/teste?x=1\#topo} pelo navegador e cole a
          requisição recebida. Por que \texttt{\#topo} não aparece, e quem
          processa esse valor?

    \item Meça as fases da conexão com um site brasileiro e outro hospedado fora
          do país. Compare DNS, TCP e TLS e proponha uma explicação para a
          diferença.
  \end{enumerate}

  \vspace{0.2cm}
  \footnotesize
  Os comandos dos itens 5 e 7 não foram vistos em aula. Estão na seção
  \textbf{Mais exemplos com curl}, no material escrito.
\end{frame}

\begin{frame}{Desafio}
  Escreva um \textbf{servidor HTTP mínimo}, em qualquer linguagem, que:

  \begin{enumerate}
    \item abra uma porta TCP;
    \item aceite uma conexão;
    \item leia a requisição recebida;
    \item devolva uma resposta HTTP válida, com \texttt{Content-Type} e
          \texttt{Content-Length} corretos.
  \end{enumerate}

  \vspace{0.4cm}
  \begin{alertblock}{Sem biblioteca pronta}
    Não use biblioteca de servidor web pronta. Trabalhe diretamente com sockets.
  \end{alertblock}

  \vspace{0.2cm}
  Em Python, cerca de vinte linhas resolvem.
\end{frame}

% ==========================================
\section*{Resumo}
% ==========================================

\begin{frame}{Resumo}
  \small
  \begin{itemize}
    \item HTTP é um acordo de formato entre cliente e servidor. O cliente sempre
          inicia.
    \item Toda mensagem tem linha inicial, cabeçalhos, linha em branco e corpo
          opcional.
    \item Métodos se distinguem por serem seguros e idempotentes. Daí decorre o
          comportamento do navegador, do cache e do histórico.
    \item Códigos de status separam sucesso, redirecionamento, culpa do cliente
          e culpa do servidor.
    \item O protocolo não guarda estado. Cookies, sessões e tokens existem para
          o cliente carregar o contexto a cada requisição.
    \item HTTPS acrescenta confidencialidade, integridade e autenticação do
          servidor, e nada mais.
    \item Origem é esquema, host e porta juntos. A restrição de mesma origem é
          do navegador, e o servidor de destino autoriza a exceção.
  \end{itemize}
\end{frame}

\begin{frame}{Bibliografia}
  \begin{itemize}
    \item TANENBAUM, A. S. \textbf{Redes de Computadores}, cap. 7.3.
    \item MDN Web Docs. \textbf{HTTP}. \\
          \url{https://developer.mozilla.org/pt-BR/docs/Web/HTTP}
    \item IETF. \textbf{RFC 9110: HTTP Semantics}, 2022. \\
          \url{https://www.rfc-editor.org/rfc/rfc9110.html}
  \end{itemize}
\end{frame}

\end{document}
